% Generated draft: requires mathematical and source review.
\documentclass[11pt]{article}
\usepackage[margin=25mm]{geometry}
\usepackage{fontspec}
\setmainfont{Latin Modern Roman}
\usepackage{amsmath,amssymb,mathtools,unicode-math}
\setmathfont{Latin Modern Math}
\usepackage{enumitem,needspace,xurl,hyperref}
\setlength{\parindent}{0pt}
\setlength{\parskip}{5pt}
\setlength{\emergencystretch}{3em}
\newcommand{\N}{\mathbb{N}}
\newcommand{\Z}{\mathbb{Z}}
\newcommand{\Q}{\mathbb{Q}}
\newcommand{\R}{\mathbb{R}}
\newcommand{\C}{\mathbb{C}}
\newcommand{\F}{\mathbb{F}}
\newcommand{\A}{\mathbb{A}}
\newcommand{\PP}{\mathbb{P}}
\newcommand{\E}{\mathbb{E}}
\newcommand{\eps}{\varepsilon}
\newcommand{\dd}{\,\mathrm{d}}
\begin{document}
\section*{UnsolvedMath ID: 10600016}
\textbf{Source problem code:} catalogue entry
\subsection{Definitions and mathematical statement}
For a diagram \(D\), let \(n(D)\) be its number of nonvirtual crossings and let \(\langle D\rangle\) be its Kauffman bracket in variable \(A\). For a nonzero Laurent polynomial, its span is the largest occurring exponent minus the smallest. Let \(g(D)\) be the source's atom (Turaev) genus. For a connected classical diagram, the usual orientable convention gives
\[
 g(D)=\frac{2+n(D)-s_A(D)-s_B(D)}2,
\]
where \(s_A,s_B\) count circles in the all-\(A\) and all-\(B\) states. The source uses the estimate
\[
 \operatorname{span}\langle D\rangle\leq4n(D)-4g(D).
\]
It reports sharpness for adequate diagrams and for many other diagrams, and recalls alternating-diagram minimality in the genus-zero setting. Adequacy means that neither all-state graph has a one-edge loop; further hypotheses such as reducedness for alternating minimality are not spelled out in the excerpt. For virtual atoms, orientability and genus conventions require specification.

The compound problem asks:
\begin{enumerate}
\item Find ways to prove that a concrete diagram realizes the minimum atom genus among diagrams of its knot. Together with equality in the displayed span estimate, such a genus-minimality result is said to imply crossing minimality in the same diagram category.
\item Explain loss of span below \(4n-4g\) in classes with minimal diagrams exhibiting that loss, including the source's torus \((p,q)\) examples with \(p>q>2\). When these are knots rather than links, coprimality is implicit and should be specified.
\item Use Kauffman brackets, Thistlethwaite spanning trees, atoms or other invariants to bound crossings of projections onto surfaces of positive genus, and establish minimality theorems in that setting. The source motivates this by drawing a torus knot without crossings on a torus. It notes that smoothing on such a surface leaves basic essential surface curves rather than only planar unknot states, so an appropriate expanded bracket target is required.
\item Determine what crossing estimates can be obtained from spanning-tree methods and/or Khovanov homology.
\item Estimate atom genus when the atom lies in a neighborhood of the projection surface rather than the plane, in the setting where the source's Khovanov thickness is at most \(2+g\).
\end{enumerate}
The genus of an individual diagram is distinguished from its minimum over all representatives. For the surface questions, the allowed surfaces, ambient embeddings, fixed versus variable genus, stabilization moves and meaning of crossing minimality are not specified. Neither the Khovanov grading/thickness convention nor the exact extended bracket is supplied.
\subsection{Short English statement}
Develop genus and crossing-minimality criteria, explain bracket-span loss, and obtain crossing and genus estimates for diagrams on higher-genus surfaces using brackets, spanning trees, atoms and Khovanov homology.
\subsection{Sources}
\begin{itemize}
\item \emph{Unsolved Problems in Virtual Knot Theory and Combinatorial Knot Theory}, revised survey, problem-list item 16, ``Crossing Number.'' \url{https://arxiv.org/abs/1409.2823}. Recorded access 2026-07-29; extraction source-tex. The statement credits Kauffman, Murasugi, Thistlethwaite and the adequate-diagram result without complete individual references.
\item The supplied summary mentions \emph{Classical results for alternating virtual links} (2022, arXiv:2204.xxxx) and Khovanov-homology work on alternating virtual links (2019, arXiv:1904.xxxx). Both identifiers are incomplete placeholders; no complete author citations are supplied.
\item Local UnsolvedMath registry, \texttt{data/unsolvedmath/problems.json}, permanent ID \(10600016\), exact source code \texttt{catalogue entry}; supplied partially\_solved/PARTIAL-PROGRESS status.
\end{itemize}
\end{document}
